\begin{figure*}[t]
\centering
\includegraphics[width=\linewidth]{figs/architecture.pdf}
\caption{\textbf{Auditable local answer path.} A prepared PDF and question enter source-scoped retrieval; selected text or pixels go to a local model; citation mapping and bounded lexical repair produce the final answer and trace. The lexical check is not an entailment judge.}
\label{fig:architecture}
\end{figure*}

\section{System Design and Trust Boundaries}
\paragraph{Canonical document state.}
Ingestion constructs a versioned bundle of source bytes, pages, structural blocks, chunks, visual regions, metadata, and content hashes. Before publication it checks source identity, pagination, evidence identifiers, image checksums, and relational constraints. Staging, per-document locking, atomic publication, and rollback prevent partially processed bundles from becoming queryable. A citation identity includes its source document, so ``Figure 2'' in one paper cannot resolve to an identically named figure in another.

\paragraph{Locality and retrieval.}
Acquisition-enabled setup may fetch documents or model assets. Strict-local analysis permits prepared uploads, caches, and loopback model services but rejects external network requests before name resolution. Query retrieval combines BM25~\citep{robertson2009bm25} and dense text rankings through reciprocal-rank fusion; crop and page-image channels provide inspection candidates for visual questions. Structural-reference expansion may bring a figure or table named by retrieved text into context. A rule-based modality heuristic can additionally boost selected visual or table candidates, but our ablation finds that this version harms first-page ranking. Retrieved pixels and model-generated visual descriptions remain different artifact types; a model description is never treated as independent source evidence.

\paragraph{One answer path and trace.}
Production and evaluation call the same versioned pipeline. It records the exact prompt evidence, raw and final completions, model and prompt identities, candidate ranks, citation origins (model-emitted versus application-added or remapped), lexical check outcomes, repair edits, abstention reason, terminal error, and stage timings. The interface exposes complete stages rather than token streaming. Figure~\ref{fig:architecture} shows the trust boundaries, not a claim that every supported-looking output is factually correct.

With those provenance boundaries fixed, the next section specifies the deliberately narrow checks and edits applied after generation.
